\textbf{Цели работы:}
\begin{itemize}
	\item  познакомиться с технологией Docker Compose;
	\item  изучить основные команды CLI;
	\item  изучить формат описания сервисов с помощью docker-compose.yml.
\end{itemize}

\textbf{Задачи:}
\begin{enumerate}
	\item изучить методические материалы к лабораторной работ;  
	\item описать и развернуть три сервиса с помощью docker-compose.yml. 
\end{enumerate}

\section*{Выполнение задания}

\begin{figure}[!htbp]
	\centering
	\includegraphics[width=0.6\linewidth]{"Screenshot 2026-03-26 140911"}
	\caption{Установленный Docker Compose}
\end{figure}
\FloatBarrier
%docker compose version

\begin{figure}[!htbp]
	\centering
	\includegraphics[width=0.6\linewidth]{"Screenshot 2026-03-26 162159"}
	\caption{Поднятые сервисы}
\end{figure}
\FloatBarrier
%docker compose ps

%https://localhost
\begin{figure}[!htbp]
	\centering
	\includegraphics[width=0.6\linewidth]{"Screenshot 2026-03-26 141048"}
	\caption{Cтраница БД}
\end{figure}
\FloatBarrier

%docker compose down -v

\section*{Исходный код программы}

\lstinputlisting[style=yml]{devops-lab4/docker-compose.yml}
Листинг 1 -- docker-compose.yml для трех сервисов
%docker compose up -d --build
%docker compose logs

\section*{Заключение}

В ходе выполнения лабораторной работы были решены следующие задачи:
\begin{itemize}
	\item  проведено ознакомление с технологией Docker Compose;
	\item  изучены основные команды CLI;
	\item  изучен формат описания сервисов с помощью docker-compose.yml.
\end{itemize}

\section*{Контрольные вопросы}

\begin{enumerate}
	\item \textit{Что такое Docker Compose и для чего он используется?}  Инструмент для определения и запуска многоконтейнерных приложений. Он упрощает управление всем стеком приложений, позволяя легко манипулировать сервисами, сетями и томами в едином и понятном конфигурационном файле YAML, а затем создать и запустить все сервисы из файла конфигурации с помощью одной команды.
	\item \textit{На рисунке ниже продемонстрирован docker-compose.yml, в котором описан один сервис db. При попытке смонтировать каталог postgres\_data внутрь контейнера для хранения записей базы данных происходит ошибка. Почему возникает ошибка и как её исправить?} 
	\begin{center}
		\includegraphics[width=0.6\linewidth]{"Screenshot 2026-03-26 135141"}
	\end{center}
	Ошибка монтирования тома, так как отсутствует объявление тома верхнего уровня. Исправляется добавлением\\volumes:\\\hspace*{2em}postgres\_data:\\\hspace*{4em}driver: local \\
	вне services.
\end{enumerate}
